\documentclass[10pt,a4paper]{amsart}
\usepackage{amsmath,amssymb,amsthm,graphicx}
\usepackage{verbatim}
\newenvironment{reference}{\comment}{\endcomment}
%\newenvironment{reference}{}{}
\newenvironment{slogan}{\comment}{\endcomment}
\newenvironment{history}{\comment}{\endcomment}

% For commutative diagrams we use Xy-pic
\usepackage[all]{xy}

% We use 2cell for 2-commutative diagrams.
\xyoption{2cell}
% Theorem environments.
%
\theoremstyle{plain}
\newtheorem{theorem}[subsection]{Theorem}
\newtheorem{proposition}[subsection]{Proposition}
\newtheorem{lemma}[subsection]{Lemma}

\theoremstyle{definition}
\newtheorem{definition}[subsection]{Definition}
\newtheorem{example}[subsection]{Example}
\newtheorem{exercise}[subsection]{Exercise}
\newtheorem{situation}[subsection]{Situation}

\theoremstyle{remark}
\newtheorem{remark}[subsection]{Remark}
\newtheorem{remarks}[subsection]{Remarks}

\numberwithin{equation}{subsection}

% Macros
%
\def\lim{\mathop{\mathrm{lim}}\nolimits}
\def\colim{\mathop{\mathrm{colim}}\nolimits}
\def\Spec{\mathop{\mathrm{Spec}}}
\def\Hom{\mathop{\mathrm{Hom}}\nolimits}
\def\Ext{\mathop{\mathrm{Ext}}\nolimits}
\def\SheafHom{\mathop{\mathcal{H}\!\mathit{om}}\nolimits}
\def\SheafExt{\mathop{\mathcal{E}\!\mathit{xt}}\nolimits}
\def\Sch{\mathit{Sch}}
\def\Mor{\mathop{\mathrm{Mor}}\nolimits}
\def\Ob{\mathop{\mathrm{Ob}}\nolimits}
\def\Sh{\mathop{\mathit{Sh}}\nolimits}
\def\NL{\mathop{N\!L}\nolimits}
\def\CH{\mathop{\mathrm{CH}}\nolimits}
\def\proetale{{pro\text{-}\acute{e}tale}}
\def\etale{{\acute{e}tale}}
\def\QCoh{\mathit{QCoh}}
\def\Ker{\mathop{\mathrm{Ker}}}
\def\Im{\mathop{\mathrm{Im}}}
\def\Coker{\mathop{\mathrm{Coker}}}
\def\Coim{\mathop{\mathrm{Coim}}}

% Boxtimes
%
\DeclareMathSymbol{\boxtimes}{\mathbin}{AMSa}{"02}

%
% Macros for moduli stacks/spaces
%
\def\QCohstack{\mathcal{QC}\!\mathit{oh}}
\def\Cohstack{\mathcal{C}\!\mathit{oh}}
\def\Spacesstack{\mathcal{S}\!\mathit{paces}}
\def\Quotfunctor{\mathrm{Quot}}
\def\Hilbfunctor{\mathrm{Hilb}}
\def\Curvesstack{\mathcal{C}\!\mathit{urves}}
\def\Polarizedstack{\mathcal{P}\!\mathit{olarized}}
\def\Complexesstack{\mathcal{C}\!\mathit{omplexes}}
% \Pic is the operator that assigns to X its picard group, usage \Pic(X)
% \Picardstack_{X/B} denotes the Picard stack of X over B
% \Picardfunctor_{X/B} denotes the Picard functor of X over B
\def\Pic{\mathop{\mathrm{Pic}}\nolimits}
\def\Picardstack{\mathcal{P}\!\mathit{ic}}
\def\Picardfunctor{\mathrm{Pic}}
\def\Deformationcategory{\mathcal{D}\!\mathit{ef}}
\usepackage{cleveref}
\crefname{originallemma}{original lemma}{original lemmas}
\Crefname{originallemma}{Original lemma}{Original lemmas}
\makeatletter
\newlabel{lemma-behaviour-of-dimensions-wrt-smooth-morphisms}{{0DRI}{source}{}{sourcepointer-1}{}}
\newlabel{lemma-behaviour-of-dimensions-wrt-smooth-morphisms@cref}{{[originallemma][0][]0DRI}{[1][0][]source}{}{}{}}
\makeatother
\crefname{originallemma}{original lemma}{original lemmas}
\Crefname{originallemma}{Original lemma}{Original lemmas}
\makeatletter
\newlabel{lemma-relative-dimension-is-semi-continuous}{{0DRQ}{source}{}{sourcepointer-2}{}}
\newlabel{lemma-relative-dimension-is-semi-continuous@cref}{{[originallemma][0][]0DRQ}{[1][0][]source}{}{}{}}
\makeatother
\crefname{originallemma}{original lemma}{original lemmas}
\Crefname{originallemma}{Original lemma}{Original lemmas}
\makeatletter
\newlabel{lemma-dimension-achieved-by-finite-type-point}{{0DRR}{source}{}{sourcepointer-3}{}}
\newlabel{lemma-dimension-achieved-by-finite-type-point@cref}{{[originallemma][0][]0DRR}{[1][0][]source}{}{}{}}
\makeatother
\title{Every point has a finite type specialization preserving dimension in any neighbourhood}
\author{The Stacks Project Authors}
\date{stacks-project:stacks-geometry.tex:local-snapshot-20261012; GNU FDL 1.2 or later}
\begin{document}
\maketitle
\noindent\textit{Editorial note.} All delivered mathematical text is the author's own, copied verbatim from the pinned Stacks Project chapter file stacks-geometry.tex; no statement or proof step was rewritten, reordered or removed.\par\medskip
\noindent\textit{Editorial note.} Span recomputation by label pairing (LANE\_BRIEF 10.2): in this chapter the author's proof environment follows the statement environment as a standalone block; the spec covers the complete environment pair (statement span, proof span) as recorded in the pool candidate.\par\medskip
\noindent\textit{Editorial note.} Occupation scan: combined claim\_key + statement\_anchor + generic-line scan over all registration files (active and withdrawn) found 0 hits for stacks-geometry.tex (ge-occupied.json, 2026-10-14); 207 route=geometry / chapter-list field hits in focused-seeds are alias noise (route name and alphabetical chapter enumerations), excluded per the ladder alias-blind-spot caveat.\par\medskip
\ifdefined\hypertarget\hypertarget{sourcepointer-1}{}\fi
\noindent\textit{Source reference.} Original lemma 0DRI: sources/stacks-chow/tags.txt, source lines 17223--17223. Cross-chapter standard prerequisite cited by the author inside the counted span; it is used as a pointer and its own proof is not part of this item.\par\medskip
\ifdefined\hypertarget\hypertarget{sourcepointer-2}{}\fi
\noindent\textit{Source reference.} Original lemma 0DRQ: sources/stacks-chow/tags.txt, source lines 17230--17230. Cross-chapter standard prerequisite cited by the author inside the counted span; it is used as a pointer and its own proof is not part of this item.\par\medskip
\ifdefined\hypertarget\hypertarget{sourcepointer-3}{}\fi
\noindent\textit{Source reference.} Original lemma 0DRR: sources/stacks-chow/tags.txt, source lines 17231--17231. Cross-chapter standard prerequisite cited by the author inside the counted span; it is used as a pointer and its own proof is not part of this item.\par\medskip
\begin{lemma}
\label{lemma-dimension-at-finite-type-point}
If $\mathcal{X}$ is a locally Noetherian algebraic stack, and if
$x \in |\mathcal{X}|$, then
for any open substack $\mathcal{V}$ of $\mathcal{X}$ containing $x$,
there is a finite type point $x_0 \in |\mathcal{V}|$ such that
$\dim_{x_0}(\mathcal{X}) = \dim_x(\mathcal{V})$.
\end{lemma}

\begin{proof}
Choose a smooth surjective
morphism  $f:U \to \mathcal{X}$ whose source is a scheme, and consider the
function $u \mapsto \dim_{f(u)}(\mathcal{X});$
since the morphism $|U| \to |\mathcal{X}|$ induced by $f$ is open (as $f$
is smooth) as well as  surjective (by assumption),
and takes finite type points to finite type points (by the very definition
of the finite type points of $|\mathcal{X}|$),
it suffices to show that for any $u \in U$, and any open neighbourhood of $u$,
there is a finite type point $u_0$ in this neighbourhood such that
$\dim_{f(u_0)}(\mathcal{X}) = \dim_{f(u)}(\mathcal{X}).$
Since, with this reformulation
of the problem,  the surjectivity of $f$ is no longer required,
we may replace $U$ by the open neighbourhood of the point $u$ in question,
and thus reduce to the problem of showing that for each $u \in U$,
there is a finite type point $u_0 \in U$ such that
$\dim_{f(u_0)}(\mathcal{X}) = \dim_{f(u)}(\mathcal{X}).$
By Lemma \ref{lemma-behaviour-of-dimensions-wrt-smooth-morphisms}
$\dim_{f(u)}(\mathcal{X}) = \dim_u(U) - \dim_u(U_{f(u)}),$
while
$\dim_{f(u_0)}(\mathcal{X}) = \dim_{u_0}(U) - \dim_{u_0}(U_{f(u_0)}).$
Since $f$ is smooth, the expression $\dim_{u_0}(U_{f(u_0)})$ is locally
constant as $u_0$ varies over $U$ (by
Lemma \ref{lemma-relative-dimension-is-semi-continuous} (2)),
and so shrinking $U$ further around
$u$ if necessary, we may assume it is constant.  Thus the problem
becomes to show that we may find a finite type point $u_0 \in U$
for which $\dim_{u_0}(U) = \dim_u(U)$.
Since by definition $\dim_u U$ is the minimum of the dimensions
$\dim V$, as $V$ ranges over the open neighbourhoods $V$ of $u$
in $U$, we may shrink $U$ down further around $u$ so that
$\dim_u U = \dim U$.
The existence of desired point $u_0$ then follows from
Lemma \ref{lemma-dimension-achieved-by-finite-type-point}.
\end{proof}

\end{document}
